\section{Partial-Order-Based Task Decomposition and Assignment}
\label{sec:po_decomp_assignment}

This section develops the first component of the top-down synthesis framework
for global temporal tasks. The starting point is the pruned task automaton
$\mathcal{B}_{\varphi}^{-}$ associated with the global sc-LTL formula
$\varphi$. The main objective is to transform the accepted automaton behavior
into a collection of executable subtask occurrences together with a compact
description of their temporal dependence. This description should be rich
enough to preserve correctness with respect to the original task, while being
light enough to support efficient assignment over a multi-robot team.

The key idea is to separate logical structure from allocation decisions.
An accepting run of $\mathcal{B}_{\varphi}^{-}$ provides a valid symbolic
execution of the global task, but the induced order is usually too restrictive.
If this linear order is used directly, then many subtasks that could be
executed concurrently would be forced into a sequential schedule. To avoid
this conservatism, the accepting run is first decomposed into indexed subtask
occurrences. The complete order of these occurrences is then relaxed into a
partial-order representation that keeps only the essential temporal
constraints. Finally, an anytime branch-and-bound procedure is applied to
assign these subtasks to feasible agents or coalitions.

The section is organized as follows.
Subsection~\ref{subsec:po_subtask_extraction} introduces the extraction of
subtask occurrences from an accepting run.
Subsection~\ref{subsec:po_relaxation} defines the partial relations and
presents the computation of accepting posets.
Subsection~\ref{subsec:po_assignment} formulates the assignment problem over
a poset and presents an anytime search algorithm.
Subsection~\ref{subsec:po_analysis} summarizes the main correctness
properties.

\subsection{Automaton-Guided Subtask Extraction}
\label{subsec:po_subtask_extraction}

Consider the pruned B\"uchi automaton
$\mathcal{B}_{\varphi}^{-}=(Q^{-},Q^{-}_{0},\Sigma,\delta^{-},Q^{-}_{F})$
obtained from the global task formula $\varphi$. Since the task is assumed to
be co-safe, satisfaction can be certified by a finite prefix that reaches an
accepting state. Therefore, the decomposition step is based on accepting runs
of finite length. Let
$\rho=q_{0}q_{1}\cdots q_{L}$ be such a run, where
$q_{0}\in Q^{-}_{0}$ and $q_{L}\in Q^{-}_{F}$. Each transition along this run
encodes one logical progress step, and thus induces one subtask occurrence.

A careful distinction must be made between a subtask label and a subtask
occurrence. The same symbolic label may appear repeatedly in an accepting run,
but those appearances can have different logical meanings because they are
attached to different automaton states. For this reason, each occurrence is
indexed explicitly. In addition, besides the enabling label that triggers the
transition to the next state, one must also retain the self-loop requirement
that characterizes the propositions allowed or required while the automaton
remains at the current state. These two components are recorded separately.

\begin{definition}[Subtask occurrence]
\label{def:po_subtask_occurrence}
Let $\rho=q_{0}q_{1}\cdots q_{L}$ be an accepting run of
$\mathcal{B}_{\varphi}^{-}$. A subtask occurrence associated with the
$\ell$th transition is defined by the triple
\begin{equation}
\label{eq:po_subtask_occurrence}
\omega_{\ell} \triangleq (\ell,\sigma_{\ell},\sigma_{\ell}^{s}),
\end{equation}
where $\ell\in\{1,\cdots,L\}$ is the occurrence index,
$\sigma_{\ell}\subseteq\Sigma$ is a minimal enabling label for the transition
from $q_{\ell-1}$ to $q_{\ell}$, and
$\sigma_{\ell}^{s}\subseteq\Sigma$ is a minimal self-loop label at
$q_{\ell-1}$. More precisely,
\begin{equation}
\label{eq:po_minimal_transition}
q_{\ell}\in\delta^{-}(q_{\ell-1},\sigma_{\ell}), \qquad
q_{\ell}\notin\delta^{-}(q_{\ell-1},\bar{\sigma}),
\end{equation}
where $\bar{\sigma}\subsetneq\sigma_{\ell}$, and
\begin{equation}
\label{eq:po_minimal_selfloop}
q_{\ell-1}\in\delta^{-}(q_{\ell-1},\sigma_{\ell}^{s}), \qquad
q_{\ell-1}\notin\delta^{-}(q_{\ell-1},\bar{\sigma}^{s}),
\end{equation}
where $\bar{\sigma}^{s}\subsetneq\sigma_{\ell}^{s}$.
The set of all occurrences induced by $\rho$ is
\begin{equation}
\label{eq:po_occurrence_set}
\Omega_{\rho} \triangleq \{\omega_{1},\cdots,\omega_{L}\}.
\end{equation}
\end{definition}

The label $\sigma_{\ell}$ in \eqref{eq:po_minimal_transition} specifies the
smallest symbolic requirement that drives the automaton from
$q_{\ell-1}$ to $q_{\ell}$. The label $\sigma_{\ell}^{s}$ in
\eqref{eq:po_minimal_selfloop} identifies the smallest symbolic requirement
that keeps the automaton at $q_{\ell-1}$. The decomposition in
\eqref{eq:po_occurrence_set} therefore captures both progress and persistence.
This distinction is important later when temporal compatibility between
subtasks is analyzed.

\begin{remark}
\label{rem:po_occurrence_index}
The occurrence index in \eqref{eq:po_subtask_occurrence} is essential.
Two subtask occurrences with the same enabling label may still need to be
treated differently if they arise at different stages of the accepting run.
Discarding the index would collapse logically distinct requests into one
symbolic object and would generally destroy the temporal structure inherited
from the automaton.
\end{remark}

The set $\Omega_{\rho}$ obtained from \eqref{eq:po_occurrence_set} is ordered
implicitly by the run index. This order is always correct, but it can be
overly conservative for execution. In particular, it forbids any concurrency
between occurrences that appear in different positions of the run, even if the
automaton would still accept after reordering some of them. The next
subsection introduces a relaxed representation that preserves only the
necessary constraints.

\subsection{Relaxed Partial-Order Abstraction}
\label{subsec:po_relaxation}

A single binary precedence relation is not sufficient for the present
setting. The reason is that action propositions are not interpreted as
instantaneous events. Instead, a proposition remains true during the entire
execution interval of the associated action. As a result, two different types
of temporal information must be represented. The first indicates that one
subtask must start no later than another. The second indicates that a group of
subtasks cannot all be active at the same time. The original source material
captures this distinction through the relations
$\preceq_{\varphi}$ and $\neq_{\varphi}$, and the same notation is adopted
here for consistency.

\begin{definition}[Partial relations]
\label{def:po_partial_relations}
Given a decomposition $\Omega_{\rho}$, define the relation
$\preceq_{\varphi}\subseteq\Omega_{\rho}\times\Omega_{\rho}$ and the family
$\neq_{\varphi}\subseteq 2^{\Omega_{\rho}}$ as follows.

For $\omega_{h},\omega_{\ell}\in\Omega_{\rho}$,
\[
\omega_{h}\preceq_{\varphi}\omega_{\ell},
\]
means that $\omega_{h}$ must be started before $\omega_{\ell}$ is started.

For a subset $X\subseteq\Omega_{\rho}$ with $|X|\geq 2$,
\[
X\in\neq_{\varphi},
\]
means that the subtasks in $X$ cannot all be executed simultaneously.
\end{definition}

The relation $\preceq_{\varphi}$ describes ordering constraints among
subtasks. The family $\neq_{\varphi}$ describes concurrent exclusion
constraints. The two notions are fundamentally different. A precedence
constraint concerns the relative order of starting times, while an exclusion
constraint concerns the overlap of execution intervals. When actions have
nonzero duration, these two notions cannot be merged without losing relevant
temporal information.

\begin{definition}[Poset of subtasks]
\label{def:po_poset}
A poset associated with the accepting run $\rho$ is the triple
\begin{equation}
\label{eq:po_poset}
P_{\rho} \triangleq
(\Omega_{\rho},\preceq_{\varphi},\neq_{\varphi}).
\end{equation}
\end{definition}

The poset in \eqref{eq:po_poset} is more general than a classical partial
order, since the concurrent exclusion information is stored by a family of
subsets instead of a binary relation. This extension is necessary because a
forbidden simultaneous execution pattern may involve more than two subtask
occurrences. The representation remains compact, however, because it records
only those relations that are logically relevant to execution.

\begin{remark}
\label{rem:po_duration}
Most automata-based planning approaches for temporal logic treat the truth of
an action proposition as instantaneous
\cite{guo2015multi,tumova2016multi,kantaros2020stylus,
luo2021abstraction,luo2022temporal,jones2019scratchs,sahin2019multirobot}.
Under that convention, exact simultaneity is rarely meaningful, and the
relation $\neq_{\varphi}$ can often be ignored. In the present setting,
however, action propositions remain valid throughout their execution
intervals. Consequently, the distinction between
$\preceq_{\varphi}$ and $\neq_{\varphi}$ is necessary.
\end{remark}

The semantic meaning of a poset is given by the set of all timed words that
respect these relations. Let $t_{\ell}$ denote the starting time of
$\omega_{\ell}$, and let $d_{\ell}$ denote its duration. The execution of
$\omega_{\ell}$ then occupies the interval $[t_{\ell},\,t_{\ell}+d_{\ell})$.

\begin{definition}[Language of a poset]
\label{def:po_language}
Given a poset
$P_{\rho}=(\Omega_{\rho},\preceq_{\varphi},\neq_{\varphi})$,
its language is the set of finite timed words
\begin{equation}
\label{eq:po_timed_word}
W_{\rho} \triangleq
(t_{1},\omega_{1})(t_{2},\omega_{2})\cdots(t_{L},\omega_{L}),
\end{equation}
such that the following two conditions hold:
\begin{equation}
\label{eq:po_prec_constraint}
t_{h}\leq t_{\ell},
\end{equation}
whenever $\omega_{h}\preceq_{\varphi}\omega_{\ell}$, and
\begin{equation}
\label{eq:po_neq_constraint}
\forall X\in\neq_{\varphi}, \;
\exists \omega_{h},\omega_{\ell}\in X
\text{ such that }
t_{h}+d_{h}<t_{\ell},
\end{equation}
where $d_{h}$ and $d_{\ell}$ are the durations of
$\omega_{h}$ and $\omega_{\ell}$, respectively.
The language of $P_{\rho}$ is denoted by $L(P_{\rho})$.
\end{definition}

Equation~\eqref{eq:po_prec_constraint} encodes the precedence requirement.
Equation~\eqref{eq:po_neq_constraint} states that, for every subset in
$\neq_{\varphi}$, at least one member must finish before another member of
that same subset starts. Hence, the entire set cannot be active
simultaneously. This interpretation follows the original formulation and is
well suited for collaborative tasks with non-negligible duration.

A poset is useful only if it is still correct with respect to the original
task. This leads to the following notion.

\begin{definition}[Accepting poset]
\label{def:po_accepting}
A poset $P_{\rho}$ is called accepting if
\begin{equation}
\label{eq:po_accepting}
L(P_{\rho}) \subseteq L_{\varphi},
\end{equation}
where $L_{\varphi}$ is the language of the global task formula.
\end{definition}

The computation starts from the fully ordered version of the accepting run and
then removes redundant relations. More specifically, let
\begin{equation}
\label{eq:po_full_order}
\mathbb{F} \triangleq
\{(\omega_{h},\omega_{\ell})\in\Omega_{\rho}\times\Omega_{\rho}
\mid h<\ell\}.
\end{equation}
The initial poset is taken as
$(\Omega_{\rho},\mathbb{F},2^{\Omega_{\rho}})$.
This choice is maximally conservative. It is then relaxed by checking
acceptance after local perturbations of the timed word associated with the
accepting run. If a swap of two adjacent occurrences still yields an accepted
word, the corresponding order relation is not essential. If a simultaneous
realization of a subset still yields an accepted word, that subset need not
remain in $\neq_{\varphi}$.

\begin{algorithm}[t!]
\caption{$\texttt{ComputePoset}(\mathcal{B}_{\varphi}^{-},t_{0})$}
\label{alg:po_extract}
\SetKwInOut{Input}{Input}
\SetKwInOut{Output}{Output}
\Input{Pruned automaton $\mathcal{B}_{\varphi}^{-}$ and time budget $t_{0}$.}
\Output{A set of accepting posets $\mathcal{P}_{\varphi}$ and their language
cover $\mathcal{L}_{\varphi}$.}

Choose $q_{0}\in Q_{0}^{-}$ and $q_{f}\in Q_{F}^{-}$\;
Initialize $\mathcal{P}_{\varphi}\leftarrow\emptyset$ and
$\mathcal{L}_{\varphi}\leftarrow\emptyset$\;
Start a modified DFS on $\mathcal{B}_{\varphi}^{-}$ to generate an accepting
run $\rho=q_{0}q_{1}\cdots q_{L}$\;

\While{$\mathrm{time}<t_{0}$}{
  Compute $\Omega_{\rho}$ from \eqref{eq:po_occurrence_set}\;
  Construct the canonical word
  $W_{\rho}\leftarrow(\omega_{1}\omega_{2}\cdots\omega_{L})$\;

  \If{$W_{\rho}\notin\mathcal{L}_{\varphi}$}{
    Set $\preceq_{\varphi}\leftarrow\mathbb{F}$ by
    \eqref{eq:po_full_order}\;
    Set $\neq_{\varphi}\leftarrow 2^{\Omega_{\rho}}$\;
    Set $P_{\rho}\leftarrow(\Omega_{\rho},\preceq_{\varphi},\neq_{\varphi})$\;
    Set $L(P_{\rho})\leftarrow\{W_{\rho}\}$ and
    $\mathrm{Que}\leftarrow\{W_{\rho}\}$\;
    Set $I_{\mathrm{acc}}\leftarrow\emptyset$ and
    $I_{\mathrm{rej}}\leftarrow\emptyset$\;

    \While{$\mathrm{Que}\neq\emptyset$}{
      Pop one word $W=(\eta_{1}\eta_{2}\cdots\eta_{L})$ from $\mathrm{Que}$\;

      \For{$i=1$ \KwTo $L-1$}{
        Set $\widetilde{W}\leftarrow
        (\eta_{1}\cdots\eta_{i+1}\eta_{i}\cdots\eta_{L})$\;

        \If{$\widetilde{W}$ is accepted by $\mathcal{B}_{\varphi}^{-}$}{
          Add $(\eta_{i},\eta_{i+1})$ to $I_{\mathrm{acc}}$\;
          Add $\widetilde{W}$ to $\mathrm{Que}$ if unseen\;
          Add $\widetilde{W}$ to $L(P_{\rho})$ if unseen\;
        }
        \Else{
          Add $(\eta_{i},\eta_{i+1})$ to $I_{\mathrm{rej}}$\;
        }
      }

      Update
      $\preceq_{\varphi}\leftarrow
      \preceq_{\varphi}\setminus
      (I_{\mathrm{acc}}\setminus I_{\mathrm{rej}})$\;

      \ForEach{$X\subseteq\Omega_{\rho}$ such that $X\in\neq_{\varphi}$}{
        Construct $\widehat{W}$ by replacing all occurrences in $X$
        within $W$ by their simultaneous union label\;

        \If{$\widehat{W}$ is accepted by $\mathcal{B}_{\varphi}^{-}$}{
          Update
          $\neq_{\varphi}\leftarrow\neq_{\varphi}\setminus\{X\}$\;
        }
      }
    }

    Update $P_{\rho}\leftarrow
    (\Omega_{\rho},\preceq_{\varphi},\neq_{\varphi})$\;
    Add $P_{\rho}$ to $\mathcal{P}_{\varphi}$\;
    Update
    $\mathcal{L}_{\varphi}\leftarrow
    \mathcal{L}_{\varphi}\cup L(P_{\rho})$\;
  }

  Continue the modified DFS until the next accepting run is found\;
}

\Return{$\mathcal{P}_{\varphi},\mathcal{L}_{\varphi}$}\;
\end{algorithm}

Algorithm~\ref{alg:po_extract} makes the relaxation process explicit.
The initial relation $\preceq_{\varphi}=\mathbb{F}$ imposes the exact order
of the accepting run, and $\neq_{\varphi}=2^{\Omega_{\rho}}$ prohibits any
simultaneous execution. The algorithm then reduces these relations by direct
acceptance tests on $\mathcal{B}_{\varphi}^{-}$. A precedence edge is removed
only if the corresponding adjacent swap is accepted. A subset is removed from
$\neq_{\varphi}$ only if its simultaneous realization is accepted. Therefore,
every relaxation is certified rather than guessed.

\begin{theorem}
\label{thm:po_accepting_poset}
Any poset returned by Algorithm~\ref{alg:po_extract} is accepting.
\end{theorem}

\begin{proof}
Consider a poset $P_{\rho}$ returned by
Algorithm~\ref{alg:po_extract}. The initial word associated with the
accepting run $\rho$ is accepted by $\mathcal{B}_{\varphi}^{-}$ by
construction. During the algorithm, a precedence relation is removed only
after the swapped word has been checked and accepted by
$\mathcal{B}_{\varphi}^{-}$. Likewise, a subset is removed from
$\neq_{\varphi}$ only after the corresponding simultaneous realization has
been checked and accepted. Hence, every timed word added to $L(P_{\rho})$
corresponds to accepted automaton behavior. It follows that
$L(P_{\rho})\subseteq L_{\varphi}$, and thus $P_{\rho}$ is accepting.
\end{proof}

\begin{remark}
\label{rem:po_compare}
The abstraction above includes decompositions based on decomposable states as
a special case \cite{schillinger2018simultaneous}. The essential difference
is that the present construction may reveal concurrency not only between
different segments of an accepting run, but also within a segment, provided
that the corresponding acceptance tests succeed. This additional flexibility
is important for minimizing the completion time of collaborative tasks.
\end{remark}

\subsection{Partial-Order-Constrained Task Assignment}
\label{subsec:po_assignment}

Once an accepting poset has been computed, the global temporal task is
reduced to an assignment problem over subtask occurrences. The logical
constraints are already encoded by $P_{\rho}$. What remains is to decide
which agents or coalitions should execute each occurrence, and in what order,
so that the resulting timed execution respects the poset and minimizes the
overall completion time.

Let $\mathcal{N}\triangleq\{1,\cdots,N\}$ be the agent index set.
For each occurrence $\omega\in\Omega_{\rho}$, define
$\Gamma(\omega)\subseteq 2^{\mathcal{N}}$ as the collection of feasible
executing coalitions. If $I\in\Gamma(\omega)$ and $|I|=1$, then $\omega$
corresponds to a local task. If $|I|>1$, then $\omega$ requires direct
collaboration among multiple agents. This notation is consistent with the
coalition-based task model introduced earlier in the chapter.

\begin{definition}[Poset graph]
\label{def:po_graph}
Given the poset
$P_{\rho}=(\Omega_{\rho},\preceq_{\varphi},\neq_{\varphi})$,
its associated poset graph is the directed graph
\begin{equation}
\label{eq:po_graph}
G_{\rho} \triangleq (\Omega_{\rho},E_{\rho}),
\end{equation}
where $(\omega_{h},\omega_{\ell})\in E_{\rho}$ if
$\omega_{h}\preceq_{\varphi}\omega_{\ell}$ and there exists no
$\omega_{m}\in\Omega_{\rho}$ such that
$\omega_{h}\preceq_{\varphi}\omega_{m}\preceq_{\varphi}\omega_{\ell}$ with
$\omega_{m}\neq\omega_{h}$ and $\omega_{m}\neq\omega_{\ell}$.
For any $\omega\in\Omega_{\rho}$, the predecessor set of $\omega$ in
$G_{\rho}$ is denoted by $\mathrm{Pred}(\omega)$.
\end{definition}

The graph in \eqref{eq:po_graph} removes all transitive edges and therefore
provides the minimal causal structure required for node expansion in the
assignment algorithm. This avoids unnecessary constraints during search while
preserving all logical dependencies encoded by $\preceq_{\varphi}$.

\begin{problem}[Poset-constrained task assignment]
\label{prob:po_assignment}
Given an accepting poset
$P_{\rho}=(\Omega_{\rho},\preceq_{\varphi},\neq_{\varphi})$,
determine for each $\omega\in\Omega_{\rho}$ a feasible coalition
$I_{\omega}\in\Gamma(\omega)$ and a start time $t_{\omega}$ such that all
constraints in \eqref{eq:po_prec_constraint} and
\eqref{eq:po_neq_constraint} are satisfied, and the makespan is minimized.
\end{problem}

Let $\mathfrak{J}(P_{\rho})$ denote the set of feasible assignments for
Problem~\ref{prob:po_assignment}. The optimal assignment is
\begin{equation}
\label{eq:po_opt_assignment}
J^{\star} \triangleq
\arg\min_{J\in\mathfrak{J}(P_{\rho})} T(J),
\end{equation}
where $T(J)$ is the completion time of assignment $J$.

Problem~\ref{prob:po_assignment} is combinatorial even without the temporal
logic structure, since it contains routing, precedence-constrained
scheduling, and coalition formation as special cases
\cite{hochba1997approximation,morrison2016branch}. For this reason, an
anytime branch-and-bound procedure is adopted. The search tree is built over
partial assignments, and only those subtasks whose predecessors have already
been assigned are eligible for expansion.

\begin{definition}[Search node]
\label{def:po_node}
A search node is a tuple
\begin{equation}
\label{eq:po_search_node}
\nu \triangleq (\tau_{1},\cdots,\tau_{N}),
\end{equation}
where $\tau_{n}$ is the ordered sequence of occurrences assigned to agent
$n\in\mathcal{N}$.
The set of assigned subtasks at node $\nu$ is
\begin{equation}
\label{eq:po_assigned_set}
\Omega_{\nu} \triangleq \bigcup_{n\in\mathcal{N}} \tau_{n}.
\end{equation}
The set of currently expandable occurrences is
\begin{equation}
\label{eq:po_expandable}
\mathrm{Next}(\nu) \triangleq
\left\{
\omega\in\Omega_{\rho}\setminus\Omega_{\nu} \;\middle|\;
\mathrm{Pred}(\omega)\subseteq\Omega_{\nu}
\right\}.
\end{equation}
\end{definition}

Equation~\eqref{eq:po_expandable} is the key admissibility condition for node
expansion. It ensures that no occurrence is assigned before all of its causal
predecessors have already appeared in the partial assignment. Hence, logical
correctness is preserved directly at the search-tree level.

The branch-and-bound method maintains one incumbent feasible assignment and
its current makespan. For every node, an upper bound is computed by greedily
completing the current partial assignment. A lower bound is computed from a
relaxation of the remaining problem. If the lower bound of a child node is no
smaller than the incumbent makespan, that branch cannot improve the current
solution and is discarded.

\begin{algorithm}[t!]
\caption{$\texttt{BnB\_Assign}(P_{\rho},t_{0})$}
\label{alg:po_bnb}
\SetKwInOut{Input}{Input}
\SetKwInOut{Output}{Output}
\Input{Accepting poset $P_{\rho}$ and time budget $t_{0}$.}
\Output{Best assignment $J^{\star}$ and best makespan $T^{\star}$.}

Initialize the root node
$\nu_{0}\leftarrow(\emptyset,\cdots,\emptyset)$\;
Initialize the active queue
$\mathcal{Q}\leftarrow\{(\nu_{0},0)\}$\;
Set $J^{\star}\leftarrow\emptyset$ and $T^{\star}\leftarrow+\infty$\;

\While{$\mathcal{Q}\neq\emptyset$ and $\mathrm{time}<t_{0}$}{
  Extract from $\mathcal{Q}$ one node $\nu$ with minimum lower bound\;

  Compute a greedy feasible completion
  $(\widehat{J}_{\nu},\widehat{T}_{\nu})$ of $\nu$\;

  \If{$\widehat{T}_{\nu}<T^{\star}$}{
    Set $J^{\star}\leftarrow\widehat{J}_{\nu}$\;
    Set $T^{\star}\leftarrow\widehat{T}_{\nu}$\;
  }

  Compute $\mathrm{Next}(\nu)$ from \eqref{eq:po_expandable}\;

  \ForEach{$\omega\in\mathrm{Next}(\nu)$}{
    \ForEach{$I\in\Gamma(\omega)$}{
      Construct the child node
      $\nu^{+}=\mathrm{Append}(\nu,\omega,I)$\;
      Compute the lower bound $\underline{T}(\nu^{+})$\;

      \If{$\underline{T}(\nu^{+})<T^{\star}$}{
        Insert $(\nu^{+},\underline{T}(\nu^{+}))$ into $\mathcal{Q}$\;
      }
    }
  }
}

\Return{$J^{\star},T^{\star}$}\;
\end{algorithm}

Algorithm~\ref{alg:po_bnb} is anytime in the usual sense.
Once the first feasible completion is found, an incumbent assignment exists.
As the search continues, the incumbent value can only improve. Therefore, the
algorithm can be terminated at any time and still returns the best feasible
solution computed so far. This property is particularly useful in large-scale
multi-robot missions, where a near-optimal solution obtained quickly may be
more valuable than an exact optimum obtained after a long delay.

\begin{theorem}
\label{thm:po_assignment_sound}
Any complete assignment returned by Algorithm~\ref{alg:po_bnb} satisfies the
poset constraints of $P_{\rho}$.
\end{theorem}

\begin{proof}
A node is expanded only with occurrences in $\mathrm{Next}(\nu)$, as defined
in \eqref{eq:po_expandable}. Therefore, every assigned occurrence appears
after all of its predecessors in the poset graph have already been assigned.
This guarantees that the precedence constraints encoded by
$\preceq_{\varphi}$ are respected. In addition, the timing synthesis step used
to evaluate a complete assignment accepts only schedules satisfying
\eqref{eq:po_neq_constraint}. Hence, every complete assignment returned by
Algorithm~\ref{alg:po_bnb} satisfies both
$\preceq_{\varphi}$ and $\neq_{\varphi}$.
\end{proof}

\subsection{Correctness and Discussion}
\label{subsec:po_analysis}

The two algorithms developed above form a coherent top-down synthesis layer.
Algorithm~\ref{alg:po_extract} computes accepting posets directly from the
task automaton. Algorithm~\ref{alg:po_bnb} searches for a low-makespan
assignment that respects one such poset. Since the first stage guarantees
logical admissibility and the second stage preserves the poset constraints,
their composition yields a sound solution method for the global temporal task.

\begin{theorem}[Soundness of the synthesis framework]
\label{thm:po_soundness}
Let $P_{\rho}$ be an accepting poset returned by
Algorithm~\ref{alg:po_extract}. Let $J^{\star}$ be a complete assignment
returned by Algorithm~\ref{alg:po_bnb} for $P_{\rho}$. Then the timed
execution induced by $J^{\star}$ satisfies the global task formula
$\varphi$.
\end{theorem}

\begin{proof}
By Theorem~\ref{thm:po_accepting_poset}, the poset $P_{\rho}$ satisfies
$L(P_{\rho})\subseteq L_{\varphi}$. By
Theorem~\ref{thm:po_assignment_sound}, the timed execution induced by
$J^{\star}$ belongs to $L(P_{\rho})$. Hence it also belongs to
$L_{\varphi}$. Therefore, the execution generated by $J^{\star}$ satisfies
the original global task formula $\varphi$.
\end{proof}

\begin{remark}
\label{rem:po_complexity}
The worst-case complexity of the overall procedure remains exponential in the
number of subtask occurrences and in the number of agents. This is expected,
since the assignment problem already contains NP-hard subproblems. The main
algorithmic advantage is not polynomial complexity, but rather the ability to
extract useful concurrency from the automaton and to maintain an incumbent
solution throughout the search. This combination is particularly attractive
for time-critical robotic applications.
\end{remark}

The present section establishes the logical and combinatorial basis of the
top-down synthesis framework. The partial-order abstraction exposes the
execution flexibility hidden in an accepting automaton run, while the
branch-and-bound assignment algorithm converts this flexibility into a concrete
multi-robot schedule. The next section can build on this representation to
address subsequent synthesis components without returning to the full global
automaton.
